% !TeX root = ../papers/nc-p3.tex
\section{縦切片 $x=1$ の円分体表示と有理水準の完全判定}
\label{sec:rank-fifteen-vertical-one}
\begin{theorem}\label{thm:rank-fifteen-vertical-one-tube}
標数零の通常な階数15窓で、剰余標数3の付値が
$\nu(y)>0$, $4\nu(x-1)\ge17\nu(y)$ を満たすとする。
窓は標準Koszul左右AS正則で、各始点で四中心四次式上自由階数256である。
特に $x=1$ では全ての正付値 $y$ に適用できる。
\end{theorem}
\begin{proof}
前節の表示で $\bar J=0$ であり、$z$ の剰余は原始八乗根である。
$\zeta^4=-1$ として $\zeta$ の上の素点を $\bar\zeta^2+\bar\zeta+2=0$
で指定し、補題\ref{lem:rank-fifteen-ramified-pole}により $z\equiv\zeta$ と選ぶ。
同じ $H$ に対して $x=1$ になる表示を次のように作る。
\[
 D=(1-2H^2)^{-1},\qquad G=\sqrt{1+2H^2D^2},\qquad
 L_0^2=\frac{D^2}{G(G+1)},\qquad C_0^2=L_0^2-1 .
\]
$G$ の定数項は1、$L_0,C_0$ の定数項は
$(\zeta+\zeta^{-1})/2,(\zeta-\zeta^{-1})/2$ とする。
これらはすべて単元で、$z_0=L_0+C_0$ は単元である。
$G^2-1=2H^2D^2$ より $G(1-2H^2L_0^2)=1$ であり、所定の縦切片となる。

\path{computations/verify_rank_fifteen_vertical_one.py} は、
$\mathbb Q[Z]/(Z^4+1)$ の係数で十七桁を取り、全ての $H$ 除算の定数項が厳密に零で
あることを検査する。単元判定には $c\equiv1\pmod3$, $c^2=-2$ を持ち上げ、
$Z^4+1=(Z^2+cZ-1)(Z^2-cZ-1)$ の第一因子を使う。
係数の共通分母の3部分が $3^e$ なら、分子を第一因子で $3^{e+1}$ まで還元する。
剰余の二係数が $3^e$ の倍数であることを確かめ、その商の少なくとも一方が
3の単元となる主元を使う。これで分母に現れる他の素点での極も扱える。
未切断経路、中心、五次の中心商に合計十桁を使う。長さ6から13の
損失列は $(1,0,1,1,2,0,1,0)$ で、一桁を残して十三次で消える。

元の $x$ については $x=G^2(1-2H^2L^2)^2$ である。
二乗の差を取ると、$1-2H^2L^2$ と $1-2H^2L_0^2$ の和、
および $L+L_0$ はいずれも単元である。
また $2(L-L_0)=(z-z_0)(1-1/(zz_0))$ の後の因子も単元なので
\[
                  \nu(x-1)=2\nu(H)+\nu(z-z_0)
\]
となる。$x=1$ なら同じ枝で $z=z_0$ である。
定理\ref{thm:rank-fifteen-collision-tube}と同じく、経路係数の変化は
$H^2(z-z_0)$ の倍数である。仮定と $\nu(y)=4\nu(H)$ から
実際の窓と円分体の表示は $H^{17}$ を法として一致する。
主元を持ち上げ、中心商の上界と一般の下界を合わせると全次数の結論を得る。
\end{proof}

\begin{corollary}\label{cor:rank-fifteen-vertical-one-rational}
標数零の代数閉体で $x=1$, $y\ne0$ とする。$y$ が代数的で原始整数最小多項式
$f$ の斉次化が $\mathbb P^1(\mathbb F_3)$ に根を持つなら、その水準の全窓が
定理の正則性と自由性を持つ。特に有理水準 $y\in\mathbb Q^\times$ の全てで成り立つ。
縦切片の添字保持同型類の座標は $t\in\mathbb G_m\setminus\{1,-1\}$ で、上記の各水準では
\[
                    y=-\frac{(t-t^{-1})^2}{108}
\]
の解がそのまま正則窓の添字保持同型類を与える。$y=1/27$ なら二類、それ以外の
上記水準では四類であり、全類の最小周期は2である。
\end{corollary}
\begin{proof}
ある3進素点で $y$ が正付値なら本定理、負付値なら
$\nu(y)<2\nu(x)$ の極の定理を使う。全ての共役根が単元の場合には、
原始式の先頭と定数はいずれも3の単元であり、斉次根は $0,\infty$ 以外である。
従って仮定から剰余根 $1$ または $-1$ があり、その上の
共役根を選べる。定理\ref{thm:rank-fifteen-unit-residues}が適用できる。
一次式の斉次化は必ず素体上に根を持つ。

この切片で $p^2=1-27y$, $q^2=-4y$ より $p^2-27q^2/4=1$ である。
$t=p+3\sqrt3\,q/2$ と取れば $p=(t+t^{-1})/2$,
$q=(t-t^{-1})/(3\sqrt3)$ となる。$q\ne0$ なので通常表示外の四点は含まれず、
定理\ref{thm:rank-fifteen-rational-quotient}と第\ref{sec:rank-fifteen-global-invariants}節の
全点分離より、相異なる解は相異なる添字保持同型類を与える。
一位置移動は $t\mapsto-t$ で固定点がなく、$p=0$ の水準だけで解は二つになる。
\end{proof}
